% Задание на выполнение ВКР (вставляется в PDF после титульника)

\newpage
\pagestyle{empty}
{%
    \begin{center}
        \footnotesize
        ПРАВИТЕЛЬСТВО РОССИЙСКОЙ ФЕДЕРАЦИИ\\[0.3em]
        ФЕДЕРАЛЬНОЕ ГОСУДАРСТВЕННОЕ АВТОНОМНОЕ ОБРАЗОВАТЕЛЬНОЕ\\
        УЧРЕЖДЕНИЕ ВЫСШЕГО ОБРАЗОВАНИЯ\\[0.3em]
        \textbf{«НАЦИОНАЛЬНЫЙ ИССЛЕДОВАТЕЛЬСКИЙ УНИВЕРСИТЕТ\\
        «ВЫСШАЯ ШКОЛА ЭКОНОМИКИ»}\\[0.8em]

        \normalsize
        Московский институт электроники и математики\\
        им.~А.Н.~Тихонова\\[2em]

        \Large\textbf{ЗАДАНИЕ}\\[0.3em]
        \large на выполнение выпускной квалификационной работы\\[1.5em]
    \end{center}

    \normalsize
    \noindent
    студенту группы~\textbf{БПМ223}\\[0.4em]
    \textbf{Бетеву Ивану Алексеевичу}\\[1em]

    \noindent\textbf{Тема работы:}\\[0.3em]
    Применение B–LoRA для повышения качества передачи стиля в
    генеративных моделях на базе DiT–FLUX.\\[1em]

    \noindent\textbf{Цель работы:}\\[0.3em]
    Исследовать применимость принципа B-LoRA к архитектуре FLUX.1 для
    управляемого разделения стиля и содержания, экспериментально оценить
    эффективность блочной адаптации относительно существующих подходов и
    установить условия её применимости.\\[1em]

    \noindent\textbf{Формулировка задания:}
    \begin{enumerate}\itemsep0.15em \setlength{\topsep}{0pt}
        \item Систематизировать существующие методы переноса стиля в
              диффузионных моделях (LoRA, IP-Adapter, SplitFlux) и выявить
              их ограничения применительно к FLUX.1.
        \item Сформулировать и экспериментально проверить гипотезу о
              стилевой специализации 19~double-stream блоков FLUX.1,
              реализующих joint attention.
        \item Разработать метод поблочного обучения стилевого
              LoRA-адаптера на double-stream блоках FLUX.1 и схему
              инференса с подключением исключительно стилевого адаптера.
        \item Сформировать экспериментальную базу: обучающие наборы стилей
              Ван~Гога и Моне, оценочные наборы на основе MS-COCO и
              ArtBench-10.
        \item Провести аблационное исследование (диапазон блоков, ранг,
              число шагов, формулировка промпта, коэффициент
              масштабирования) и выбрать оптимальную конфигурацию метода.
        \item Выполнить сравнение предложенного метода с Full-LoRA-FLUX,
              SplitFlux и оригинальным B-LoRA-SDXL по метрикам
              DINO ViT-B/8, CLIP-style, CLIP-content, FID, LPIPS.
        \item Проанализировать режимы отказа метода и сформулировать
              выводы об условиях его применимости.
    \end{enumerate}

    \vspace{1.2em}
    \noindent Проект ВКР должен быть предоставлен студентом в срок до\\[0.2em]
    \mbox{«\,\rule{2em}{0.4pt}\,»\,\rule{8em}{0.4pt}\,20\rule{1.5em}{0.4pt}\,г.}\\[0.4em]
    Научный руководитель ВКР\\[0.2em]
    \mbox{«\,\rule{2em}{0.4pt}\,»\,\rule{8em}{0.4pt}\,20\rule{1.5em}{0.4pt}\,г.}\hfill
    \rule{8em}{0.4pt}\quad С.А.~Сластников\\[0.8em]

    \noindent Первый вариант ВКР предоставлен студентом в срок до\\[0.2em]
    \mbox{«\,\rule{2em}{0.4pt}\,»\,\rule{8em}{0.4pt}\,20\rule{1.5em}{0.4pt}\,г.}\\[0.4em]
    Научный руководитель ВКР\\[0.2em]
    \mbox{«\,\rule{2em}{0.4pt}\,»\,\rule{8em}{0.4pt}\,20\rule{1.5em}{0.4pt}\,г.}\hfill
    \rule{8em}{0.4pt}\quad С.А.~Сластников\\[0.8em]

    \noindent Итоговый вариант ВКР предоставлен студентом в срок до\\[0.2em]
    \mbox{«\,\rule{2em}{0.4pt}\,»\,\rule{8em}{0.4pt}\,20\rule{1.5em}{0.4pt}\,г.}\\[0.4em]
    Научный руководитель ВКР\\[0.2em]
    \mbox{«\,\rule{2em}{0.4pt}\,»\,\rule{8em}{0.4pt}\,20\rule{1.5em}{0.4pt}\,г.}\hfill
    \rule{8em}{0.4pt}\quad С.А.~Сластников\\[1.5em]

    \noindent Задание выдано студенту\\[0.2em]
    \mbox{«\,\rule{2em}{0.4pt}\,»\,\rule{8em}{0.4pt}\,20\rule{1.5em}{0.4pt}\,г.}\hfill
    \rule{8em}{0.4pt}\quad С.А.~Сластников\\
    \hspace*{0pt}\hfill\textit{(научный руководитель ВКР)}\hspace*{1em}\\[0.8em]

    \noindent Задание принято к исполнению студентом\\[0.2em]
    \mbox{«\,\rule{2em}{0.4pt}\,»\,\rule{8em}{0.4pt}\,20\rule{1.5em}{0.4pt}\,г.}\hfill
    \rule{8em}{0.4pt}\quad И.А.~Бетев\\
}
\clearpage
\pagestyle{plain}
